\chapter{To Have, and How Old You Are}
\label{ch:haben-alter}

German's workhorse verb \emph{haben}, ``to have'' --- which hides the best false
cognate in the whole course --- and then the one everyday sentence where German
\emph{refuses} to use it, siding with English against every Romance sister.
\emph{Practice lessons: \texttt{lessons/GE-C14-*.md}.}

\section{\emph{haben} --- ``to have,'' the twin of English \emph{have}}
\label{lesson:haben}

\begin{sounds}
\emph{HAH-ben}, with a long ``a'' as in English \emph{father}. In \emph{hast}
and \emph{hat} the vowel stays short and the final \emph{t} is crisp.
\end{sounds}

\begin{center}
\begin{tabular}{@{}ll@{}}
\toprule
Singular & Plural \\
\midrule
ich \textbf{habe} & wir \textbf{haben} \\
du \textbf{hast} & ihr \textbf{habt} \\
er/sie \textbf{hat} & sie/Sie \textbf{haben} \\
\bottomrule
\end{tabular}
\end{center}

\begin{morphologybox}
Mostly regular, with \textbf{two irregulars}: \emph{du hast} and \emph{er hat}
\textbf{drop the b}. English does the identical thing --- \emph{have} but
\emph{ha\textbf{s}}, and \emph{hast} in older English. Two languages, one
shortcut, inherited together.
\end{morphologybox}

\subsection*{The trap: \emph{haben} is not \emph{hab\=ere}}

\begin{center}
\begin{tabular}{@{}lll@{}}
\toprule
Language & Verb & Meaning \\
\midrule
German & \textbf{haben} & to have \\
Latin & \textbf{hab\=ere} & to have \\
\bottomrule
\end{tabular}
\end{center}

Same consonants, same meaning. \textbf{They are not related.} It is a
coincidence --- one of the most famous in linguistics.

\begin{cousinweb}
\textbf{haben} $\leftarrow$ Germanic \emph{*habjan\k{a}} $\leftarrow$
Proto-Indo-European \emph{*kap-}, ``\textbf{to seize}.'' That root's Latin child
is \emph{capere} --- which gives English \textbf{capture}, \textbf{captive},
\textbf{capable}, \textbf{accept}.\\[3pt]
\textbf{hab\=ere} $\leftarrow$ Proto-Indo-European \emph{*g\ensuremath{^h}ab\ensuremath{^h}-},
whose English descendant is\ldots{} \textbf{give}.\\[3pt]
So German \emph{haben} is a cousin of \textbf{capture}, and Latin \emph{hab\=ere}
--- the source of French \emph{avoir} and Italian \emph{avere} --- is a cousin of
\textbf{give}. The two words that look most alike come from \textbf{opposite}
directions.
\end{cousinweb}

\subsection*{Using it}

German nouns are capitalised and take their gender:

\begin{center}\itshape
Ich habe einen Bruder. --- ``I have a brother.''\\
Wir haben Brot und Wein. --- ``We have bread and wine.''
\end{center}

Say the whole set aloud --- \emph{ich habe, du hast, er hat, wir haben, ihr
habt, sie haben} --- listening for the dropped \textbf{b}: \emph{habe \ldots{}
ha\textbf{st} \ldots{} ha\textbf{t}}.

\section{\emph{ich bin zwanzig Jahre alt} --- being your years}
\label{lesson:alter}

You have just learned \emph{haben}. Here is the everyday sentence where German
\textbf{won't} let you use it --- and where German and English stand alone
against all four Romance languages.

\begin{center}\itshape
Wie alt bist du? --- literally ``How old are you?''\\
Ich bin zwanzig Jahre alt. --- ``I am twenty years old.''
\end{center}

Word for word, that is the \textbf{English sentence}. Not a coincidence --- both
languages inherited the same way of saying it. You can shorten it to \emph{Ich
bin zwanzig}, exactly as English drops ``years old.''

\begin{cousinweb}
\textbf{Jahr} $\leftarrow$ Germanic \emph{*j\=era} $=$ English \textbf{year}.\\[3pt]
\textbf{alt} $\leftarrow$ Germanic \emph{*aldaz} $=$ English \textbf{old}.\\[3pt]
\emph{Alt} has a Latin cousin too: \emph{alere}, ``to nourish, grow'' --- which
gives English \textbf{adult} (``grown up'') and \textbf{old} itself. Growing and
ageing, one idea.
\end{cousinweb}

\begin{morphologybox}
\textbf{Jahre} is the plural of \emph{Jahr}, and German capitalises it because
it is a noun. The same word sits inside \emph{Jahreszeiten} --- ``year's-times,''
the German word for the seasons.
\end{morphologybox}

\subsection*{The split, laid out}

\begin{center}
\begin{tabular}{@{}lll@{}}
\toprule
Language & How age works & Literally \\
\midrule
\textbf{German} & \emph{ich \textbf{bin} \ldots{} Jahre alt} & I \textbf{am} \ldots{} \\
\textbf{English} & \emph{I \textbf{am} twenty} & I \textbf{am} \ldots{} \\
French & \emph{j'\textbf{ai} vingt ans} & I \textbf{have} \ldots{} \\
Spanish & \emph{t\textbf{engo} veinte a\~{n}os} & I \textbf{have} \ldots{} \\
Italian & \emph{h\textbf{o} venti anni} & I \textbf{have} \ldots{} \\
Portuguese & \emph{t\textbf{enho} vinte anos} & I \textbf{have} \ldots{} \\
\bottomrule
\end{tabular}
\end{center}

\begin{culture}
\textbf{Romance has its years; Germanic is its years.} Age as something you
carry versus age as something you are --- and German, which borrowed its month
names from Latin, keeps its \textbf{own} logic here. The verb it uses is
\emph{sein} (\emph{bin}), not \emph{haben}; that verb gets a chapter of its own
shortly.
\end{culture}
